\section{Order filters, neutral curves and the digamma inverse}
\label{ojf:section}

The preceding sections follow singularities through feedback and inversion.
We now ask a complementary design question: which exponentially small data
can an order derivative cancel before inversion is performed? For the
Fermi family, exact completion turns this into interpolation at conjugate
logarithmic nodes. The resulting filters lead to moving-order curves with
zero lateral difference, whose physical inverses still have a nonzero
exponential displacement. At order zero, the same completion gives a
digamma core and a globally invertible positive integral.

This develops the order-filter continuation \cite{ojf:report} of the
resonant polylogarithm report \cite{qfb:reports}. Its Fermi limit
$-e^x\to-\infty$ differs from the fixed-sheet Lerch limit toward one
earlier in this chapter. The common tools are order differentiation,
exact completion and controlled local inversion. The neutral inverse
chart below resolves a further question left open in that continuation.

A terminating or factorially divergent power--logarithmic expansion does
not determine a unique analytic function. Here the distinction
can be made exact. For $x>0$ set
\begin{equation}\label{ojf:eq:defs}
 F_s(x)=-\Li_s(-e^x),\qquad
 T_s(x)=F_s(-x)=\sum_{n\ge1}\frac{(-1)^{n-1}e^{-nx}}{n^s}.
\end{equation}
The subscript is the \emph{polylogarithm order}; the usual normalized
Fermi--Dirac index is $s-1$ \cite{rw:DLMFPolylog}.
There are two distinguished lateral realizations $L_s^\pm$ of the
Sommerfeld expansion, satisfying
\begin{equation}\label{ojf:eq:intro-completion}
 L_s^\pm=F_s+e^{\mp\iu\pi s}T_s.
\end{equation}
Consequently the jump $L_s^+-L_s^-$ and the physical correction
$(L_s^++L_s^-)/2-F_s$ are different data. Neither should be silently
substituted for the other. We use \emph{arithmetic average}, not
``median,'' for $(L_s^++L_s^-)/2$.

For a real polynomial $\Pi$, applying $\Pi(\partial_s)$ changes the
coefficient of the $n$th exponential into the value of $\Pi$ at
$-\log n\mp\iu\pi$. We ask how many initial actions can be removed by
an order jet of prescribed degree, how this cancellation behaves under
small order detuning, and what analytic inversion does to the remaining
completion. Here an \emph{action $n$} simply means the exponential
$e^{-nx}$; this terminology does not add a summability hypothesis.

\subsection{The Fermi family and exact lateral completion}
\label{ojf:sec:completion}

\subsubsection{The analytic realization used throughout}
We first state the existing completion in the exact normalization needed
below, and give its derivation to fix signs and domains.

\begin{proposition}[Classical completion in the repository normalization]
\label{ojf:prop:completion}
For $x>0$ and $\Real s>0$, the formal Sommerfeld series is
\begin{equation}\label{ojf:eq:sommerfeld}
 \widehat P_s(x)=\frac{x^s}{\Gamma(s+1)}+
       2\sum_{k\ge1}\frac{\eta(2k)}{\Gamma(s-2k+1)}x^{s-2k},
 \qquad \eta(z)=(1-2^{1-z})\zeta(z).
\end{equation}
In the Borel convention
$\mathcal{B}(\sum_{j\ge1}a_jx^{-j})(\xi)=
 \sum_{j\ge1}a_j\xi^{j-1}/(j-1)!$,
the normalized minor is
\begin{equation}\label{ojf:eq:lattice}
 b_s(\xi)=s\sum_{n\ge1}\frac{(-1)^{n-1}}n
 \left\{(1+\xi/n)^{s-1}-(1-\xi/n)^{s-1}\right\}.
\end{equation}
The upper and lower lateral realizations
\begin{equation}\label{ojf:eq:lateral-borel}
 L_s^\pm(x)=\frac{x^s}{\Gamma(s+1)}
       \left(1+\int_{0}^{\infty e^{\pm\iu0}}
                      e^{-x\xi}b_s(\xi)\dd\xi\right)
\end{equation}
satisfy \eqref{ojf:eq:intro-completion}. Equivalently, with
$a_\pm=\frac12\pm\frac{\iu x}{2\pi}$,
\begin{equation}\label{ojf:eq:hurwitz-completion}
 L_s^\pm(x)=-\frac{(2\pi)^se^{\mp\iu\pi s/2}}{\Gamma(s)}
                     \zeta(1-s,a_\pm).
\end{equation}
The last formula has a removable singularity at $s=0$ and is entire in
$s$. It supplies the continuation of the lateral realization to every
order. All finite order derivatives may be taken in the exact completion.
\end{proposition}

\begin{proof}
For $\Real s>0$ the Fermi integral is
\[
 F_s(x)=\frac1{\Gamma(s)}\int_0^\infty
                    \frac{t^{s-1}}{1+e^{t-x}}\dd t.
\]
Split at $t=x$ and put $u=|t-x|$. This gives
\begin{equation}\label{ojf:eq:thermal-split}
 F_s(x)=\frac{x^s}{\Gamma(s+1)}+
 \frac1{\Gamma(s)}\int_0^\infty
 \frac{(x+u)^{s-1}-{\bf1}_{u<x}(x-u)^{s-1}}{e^u+1}\dd u.
\end{equation}
The binomial expansion of \eqref{ojf:eq:lattice} at zero has coefficient
$2\eta(2k)s^{\underline{2k}}/(2k-1)!$ at $\xi^{2k-1}$.
This proves the formal Borel identity. The lattice sum converges normally
on compact continued sheets away from $\pm1,\pm2,\ldots$: the bracket
in its $n$th term is $O(n^{-1})$ and its prefactor contributes another
$n^{-1}$. The same estimates, with logarithmic factors, hold for finitely
many order derivatives on compact order sets.

On a ray strictly above or below the positive axis the kernel has at most
polynomial--logarithmic growth. One can verify this by splitting the sum
at $n\asymp|\xi|$: the larger $n$ use the preceding difference estimate,
and the finitely many smaller $n$ use the power growth of their binomials.
Thus the tilted Laplace integrals converge. On the boundary their local
powers are integrable when $\Real s>0$.

In the $n$th integral substitute $\xi=nu/x$. Summing
$\sum_{n\ge1}(-1)^{n-1}e^{-nu}=1/(e^u+1)$ turns the integral part
of \eqref{ojf:eq:lateral-borel} into
\[
 \frac1{\Gamma(s)}\int_0^\infty
 \frac{(x+u)^{s-1}-(x-u)_{\mp}^{s-1}}{e^u+1}\dd u.
\]
The upper path gives argument $-\pi$ to $x-u$ when $u>x$.
Subtracting \eqref{ojf:eq:thermal-split} therefore gives
\[
 -\frac{e^{\mp\iu\pi(s-1)}}{\Gamma(s)}
  \int_x^\infty\frac{(u-x)^{s-1}}{e^u+1}\dd u
   =e^{\mp\iu\pi s}T_s(x).
\]
Endpoint domination by $v^{\epsilon-1}(1+|\log v|^r)$ on compact
sets $\Real s\ge\epsilon>0$, and exponential domination at infinity,
justify the differentiations and interchanges. Finite Taylor expansion
of the Borel germ with a remainder gives the asymptotic interpretation
of \eqref{ojf:eq:sommerfeld}; no infinite divergent sum was integrated
term by term.

Jonqui\`ere's identity \cite[Eq.~25.12.13]{rw:DLMFPolylog}, at $z=-e^x$,
gives the lower formula in \eqref{ojf:eq:hurwitz-completion}; conjugating
for real $s$ and then using analytic continuation gives the upper one.
For $a_\pm$ in the right half-plane, $\zeta(1-s,a_\pm)$ has only its
simple pole at $s=0$, cancelled by the simple zero of $1/\Gamma(s)$.
The small-argument series $T_s$ is entire in $s$, normally on compact
sets. Hence both the completion and the physical $F_s$ continue
entirely in the order, and finite order derivatives remain valid.
\end{proof}

\begin{remark}[Continuation is not an unregularized boundary integral]
For $\Real s\le0$, \eqref{ojf:eq:hurwitz-completion} is the definition of
the continued lateral realization used here. We do not assert that the
boundary powers in \eqref{ojf:eq:lateral-borel} become ordinary locally
integrable functions there. In Subsection~\ref{ojf:sec:borelzero} we give a
separate, ordinary meromorphic Borel representation for the first jet
at zero. Differentiating the \emph{whole} normalized expression matters:
$\log x$ is not treated as a constant coefficient in an unjustified
ordinary Borel transform.
\end{remark}

\subsubsection{A polynomial dictionary for the exponential data}
For real $\sigma$ and $\Pi\in\R[u]$, write
\begin{align}
 \mathcal{F}_{\sigma,\Pi}(x)&=
       \left.\Pi(\partial_s)F_s(x)\right|_{s=\sigma},\label{ojf:eq:observable}\\
 \mathcal{L}^\pm_{\sigma,\Pi}(x)&=
       \left.\Pi(\partial_s)L_s^\pm(x)\right|_{s=\sigma},\notag\\
 C_{\sigma,\Pi}&=\tfrac12(\mathcal{L}^+_{\sigma,\Pi}+
                           \mathcal{L}^-_{\sigma,\Pi})-\mathcal{F}_{\sigma,\Pi},\notag\\
 J_{\sigma,\Pi}&=\mathcal{L}^+_{\sigma,\Pi}-\mathcal{L}^-_{\sigma,\Pi}.\notag
\end{align}
Put $D=\mathrm d/\mathrm du$. The operators
\begin{equation}\label{ojf:eq:operators}
 R_\sigma=\cos(\pi\sigma+\pi D),\qquad
 I_\sigma=\sin(\pi\sigma+\pi D)
\end{equation}
are their Taylor series in the nilpotent operator $D$ on any fixed-degree
polynomial space; all sums terminate.

\begin{theorem}[Exact order-filter dictionary]\label{ojf:thm:dictionary}
Let $u_n=-\log n$ and $q=e^{-x}$. Then
\begin{align}
 \mathcal{L}^\pm_{\sigma,\Pi}-\mathcal{F}_{\sigma,\Pi}
 &=e^{\mp\iu\pi\sigma}
       \sum_{n\ge1}\frac{(-1)^{n-1}}{n^\sigma}
                      \Pi(u_n\mp\iu\pi)q^n,\label{ojf:eq:filter-full}\\
 C_{\sigma,\Pi}
 &=\sum_{n\ge1}\frac{(-1)^{n-1}}{n^\sigma}(R_\sigma\Pi)(u_n)q^n,
                                                        \label{ojf:eq:filter-C}\\
 J_{\sigma,\Pi}
 &=-2\iu\sum_{n\ge1}\frac{(-1)^{n-1}}{n^\sigma}(I_\sigma\Pi)(u_n)q^n.
                                                        \label{ojf:eq:filter-J}
\end{align}
The series converge absolutely for $|q|<1$, normally on compact order
sets. For real $u$,
\begin{equation}\label{ojf:eq:rotation}
 (R_\sigma\Pi)(u)+\iu(I_\sigma\Pi)(u)
      =e^{\iu\pi\sigma}\Pi(u+\iu\pi).
\end{equation}
\end{theorem}
\begin{proof}
Every derivative of
$e^{\mp\iu\pi s}n^{-s}$ multiplies it by
$-\log n\mp\iu\pi$. Apply the finite operator to
\eqref{ojf:eq:intro-completion}. A polynomial in $\log n$ times
$n^{-\sigma}$ has subexponential growth in $n$, which proves the stated
normal convergence. Taylor's formula for a polynomial gives
$\Pi(u+\iu\pi)=\exp(\iu\pi D)\Pi(u)$ and hence
\eqref{ojf:eq:rotation}. Taking the arithmetic average and difference
proves the two remaining formulas, with the sign in the jump fixed by
upper minus lower.
\end{proof}

\subsection{Designing order filters: separate and joint cancellation}
\label{ojf:sec:filters}

\subsubsection{What a degree can and cannot remove}
For a nonzero power series in $q$ with no constant term, let $\nu$ be its
first nonzero exponent; let $\nu(0)=\infty$.

\begin{theorem}[Sharp separate degree costs]\label{ojf:thm:cost}
Fix $\Pi\ne0$ of degree $d$ and real $\sigma$.
If $\sin\pi\sigma\ne0$, then $\nu(J_{\sigma,\Pi})\le d+1$;
if $\sin\pi\sigma=0$ and $d\ge1$, then
$\nu(J_{\sigma,\Pi})\le d$.
Similarly, if $\cos\pi\sigma\ne0$, then
$\nu(C_{\sigma,\Pi})\le d+1$;
if $\cos\pi\sigma=0$ and $d\ge1$, then
$\nu(C_{\sigma,\Pi})\le d$.
Every displayed bound is attained by a polynomial of the indicated degree.
The exceptional constant polynomials have zero jump at integer order,
and zero arithmetic completion at half-integer order.
\end{theorem}

\begin{proof}
The constant term of $I_\sigma$ as a polynomial in $D$ is
$\sin\pi\sigma$. When this is nonzero, $I_\sigma$ is an invertible
triangular operator on degree-at-most-$d$ polynomials and preserves the
degree. When it is zero,
\[
 I_\sigma=(\pi\cos\pi\sigma)D+O(D^3),
\]
so it has kernel exactly the constants and reduces every positive degree
by exactly one. The analogous assertions for $R_\sigma$ follow from
its constant and linear terms, $\cos\pi\sigma$ and
$-\pi\sin\pi\sigma$.

A nonzero polynomial of degree $k$ cannot vanish at all of the $k+1$
distinct real numbers $u_1,\ldots,u_{k+1}$. The coefficients multiplying
these evaluations in \eqref{ojf:eq:filter-C}--\eqref{ojf:eq:filter-J}
never vanish. This proves the upper bounds.

To attain a bound for a degree-preserving operator, take its inverse
image of $\prod_{n=1}^d(u+\log n)$. For a degree-lowering operator,
its map from polynomials of degree at most $d$, modulo constants, to
polynomials of degree at most $d-1$ is an isomorphism. Take the inverse
image of $\prod_{n=1}^{d-1}(u+\log n)$. The next node is not a root,
so the bound is attained exactly. The assertions for constants follow
directly from \eqref{ojf:eq:operators}.
\end{proof}

Thus, for $d\ge1$, the optimal first surviving actions are
\begin{center}
\begin{tabular}{@{}lll@{}}\toprule
Order center & Jump & Arithmetic completion\\\midrule
$\sigma\in\mathbb Z$ & $d$ & $d+1$\\
$\sigma\in\mathbb Z+\tfrac12$ & $d+1$ & $d$\\
Neither of the above & $d+1$ & $d+1$\\\bottomrule
\end{tabular}
\end{center}
The two columns are separately optimized; they need not be realized by
the same polynomial. In particular, a degree statement about one datum
cannot be reused for their simultaneous cancellation.

For example, at integer order the operators are, up to $(-1)^\sigma$,
$\cos(\pi D)$ and $\sin(\pi D)$. A polynomial that delays the completion
to action $d+1$ is
\begin{equation}\label{ojf:eq:cos-extremizer}
 \Pi^{C}_d(u)=\sec(\pi D)\prod_{n=1}^d(u+\log n).
\end{equation}
The power series for $\sec$ terminates on the polynomial. A primitive
followed by the finite operator $\pi D/\sin(\pi D)$ similarly constructs
a jump extremizer. These formulas are exact algebra, not asymptotic
infinite operator expansions.

\subsubsection{Simultaneous cancellation is a principal ideal}

\begin{theorem}[Joint cancellation ideal and optimal real jet degree]
\label{ojf:thm:joint}
Let $E$ be any finite set of positive integers. For every real order
center $\sigma$, the real polynomials $\Pi$ for which both the jump and
the arithmetic completion vanish at every action in $E$ form exactly
\begin{equation}\label{ojf:eq:joint-ideal}
 \left(\prod_{n\in E}\bigl((u+\log n)^2+\pi^2\bigr)\right)\R[u].
\end{equation}
Consequently, for $E=\{1,\ldots,N\}$ the minimum nonzero degree is $2N$,
and the unique monic minimizer is
\begin{equation}\label{ojf:eq:joint-extremizer}
 Q_N^{\mathrm{joint}}(u)=
       \prod_{n=1}^N\bigl((u+\log n)^2+\pi^2\bigr).
\end{equation}
At action $N+1$, at least one of its jump and completion coefficients
is nonzero. More generally a degree-$d$ nonzero jet cannot make both
data zero at all of the first $\lfloor d/2\rfloor+1$ actions.
\end{theorem}

\begin{proof}
By \eqref{ojf:eq:rotation}, both real components at the node $u_n$
vanish if and only if $\Pi(u_n+\iu\pi)=0$. Since $\Pi$ is real,
this is equivalent to zeros at both $u_n+\iu\pi$ and $u_n-\iu\pi$,
or to divisibility by $(u-u_n)^2+\pi^2$. These real quadratic factors
are pairwise coprime: all their complex roots are distinct. Divisibility
by each is therefore equivalent to divisibility by their product.
This proves the ideal identity and the degree assertions.

For the monic minimizer at the next node, putting
$\ell_j=\log(j/(N+1))$ gives
\begin{equation}\label{ojf:eq:next-joint}
 Q_N^{\mathrm{joint}}(-\log(N+1)+\iu\pi)
      =\prod_{j=1}^N\ell_j(\ell_j+2\iu\pi)\ne0.
\end{equation}
After rotation by $e^{\iu\pi\sigma}$ this still has at least one
nonzero real component. This proves the final assertion. We do not
claim that both components must separately be nonzero in this formula.
\end{proof}

\begin{corollary}[First-order detuning robustness]\label{ojf:cor:robustness}
For a fixed real polynomial $\Pi$ and a finite action set $E$, the
following requirements are equivalent: the jump coefficients at every
$n\in E$ vanish together with their first order-center derivatives at
one real center $\sigma_0$; both the jump and completion coefficients
vanish there; and the full lateral completion coefficients vanish at
those actions for \emph{every} order center. Their common solution ideal
is \eqref{ojf:eq:joint-ideal}. In particular robustness through first
order in detuning costs exactly $2N$ for the first $N$ actions.
\end{corollary}
\begin{proof}
Up to a nonzero fixed scalar, the $n$th jump coefficient is
$n^{-\sigma}(I_\sigma\Pi)(u_n)$. Its derivative is
\[
 n^{-\sigma}\{-\log n\,(I_\sigma\Pi)(u_n)
                      +\pi(R_\sigma\Pi)(u_n)\}.
\]
When the jump coefficient is zero, its derivative is zero exactly when
the completion coefficient is zero. The joint theorem then forces
$\Pi(u_n\pm\iu\pi)=0$, independently of the order center. The converse
is immediate from the exact dictionary.
\end{proof}

\begin{corollary}[Finite reconstruction from exponential data]
\label{ojf:cor:reconstruction}
Fix $N$ and real $\sigma$. A polynomial of degree less than $2N$ is
uniquely determined by its jump and arithmetic-completion coefficients
at the first $N$ actions. Every assignment of those $2N$ real normalized
data is attained by exactly one such polynomial.
\end{corollary}
\begin{proof}
Divide out the nonzero scalar coefficients in
\eqref{ojf:eq:filter-C}--\eqref{ojf:eq:filter-J} and undo the rotation
in \eqref{ojf:eq:rotation}. This specifies $\Pi$ at the $N$ upper nodes
and, by conjugacy, at the $N$ lower nodes. Ordinary interpolation at
$2N$ distinct complex points gives a unique degree-less-than-$2N$
polynomial. Its conjugate-reflection polynomial has the same values,
so uniqueness makes its coefficients real. Equivalently, injectivity
follows from Theorem~\ref{ojf:thm:joint} and equal finite dimensions give
surjectivity.
\end{proof}

\paragraph{The first joint filter.}
For $N=1$, the operator is $\partial_s^2+\pi^2$. It removes the
entire first action, not just its imaginary part. At integer $m$ the
next completion coefficient is
$(-1)^{m+1}(\log2)^2/2^m$ and the next jump coefficient is
$-4\pi\iu(-1)^m\log2/2^m$. This is a simple exact test of the
real-versus-imaginary distinction.

\subsubsection{An explicit finite-tail estimate}
For numerical evaluation the exact coefficient dictionary also supplies
an elementary bound. Let $\Pi(u)=\sum_{j=0}^dp_ju^j$ and
$K_\Pi=\sum_{j=0}^d|p_j|(1+\pi)^j$. For real $\sigma\ge0$,
$0<q<1$, and integer $M\ge0$, each full lateral completion tail obeys
\begin{equation}\label{ojf:eq:tailbound}
 \left|\sum_{n>M}\frac{(-1)^{n-1}}{n^\sigma}
                  \Pi(-\log n\pm\iu\pi)q^n\right|
 \le K_\Pi(M+1)^{d-\sigma}q^{M+1}
                  \frac{E_d(q)}{(1-q)^{d+1}},
\end{equation}
where $E_d$ is the polynomial defined by
$\sum_{k\ge0}(k+1)^dq^k=E_d(q)/(1-q)^{d+1}$, with $E_0=1$.
Indeed $|\Pi(-\log n\pm\iu\pi)|\le K_\Pi n^d$ for $n\ge1$,
$n=M+1+k\le(M+1)(k+1)$, and $n^{-\sigma}\le(M+1)^{-\sigma}$.
This proves the bound. The jump tail has at most twice the right side;
the arithmetic-completion tail has at most the right side. This is a
finite evaluation safeguard, not the basis of the cancellation theorem.

\subsection{Arithmetic rigidity and moving-order cancellation}
\label{ojf:sec:neutral}

\subsubsection{A transcendence input rules out accidental zeros}

\begin{lemma}\label{ojf:lem:log-transcendence}
For every integer $n\ge2$, $\log n/\pi$ is transcendental.
\end{lemma}
\begin{proof}
Use the classical Gelfond--Schneider theorem in its complex form: if
$\alpha\ne0,1$ and $\beta\notin\Q$ are algebraic, every fixed value of
$\alpha^\beta$ is transcendental \cite[Theorem 3.1]{ojf:GS}.
If $a=\log n/\pi$ were algebraic, $\beta=-\iu a$ would be nonreal
algebraic, hence irrational. With $\log(-1)=\iu\pi$,
$(-1)^{-\iu a}=\exp(\pi a)=n$, a contradiction.
\end{proof}

\begin{theorem}[Parity and nonvanishing for every pure jet]
\label{ojf:thm:parity}
Fix an integer order $m$ and $r\ge1$, and take $\Pi(u)=u^r$.
At every action $n\ge2$, both the jump coefficient and the
arithmetic-completion coefficient are nonzero. Their first surviving
actions are exactly
\begin{equation}\label{ojf:eq:parity-table}
 (\nu(J_{m,u^r}),\nu(C_{m,u^r}))=
 \begin{cases}(1,2),&r\ \text{odd},\\(2,1),&r\ \text{even}.
 \end{cases}
\end{equation}
For $r=0$, the jump is identically zero and the completion starts at
one. At half-integer centers the two columns of
\eqref{ojf:eq:parity-table} are exchanged.
\end{theorem}
\begin{proof}
At $n=1$, the relevant complex value is $(\iu\pi)^r$; its real and
imaginary parts alternate with the parity of $r$. At $n\ge2$, divide
$(-\log n+\iu\pi)^r$ by $\pi^r$. Its real and imaginary parts are
nonzero polynomials in $a=\log n/\pi$ with integer coefficients:
the real polynomial has leading term $(-a)^r$, and the imaginary one
has leading term $r(-a)^{r-1}$. Neither can vanish at the transcendental
number in Lemma~\ref{ojf:lem:log-transcendence}. At integer order the rotation
is just $(-1)^m$; at half-integer order it is $\pm\iu$, which exchanges
the two components. The constant case follows from the completion formula.
\end{proof}

\begin{corollary}[Rigidity for algebraic normalized coefficients]
\label{ojf:cor:algebraic-rigidity}
The nonvanishing at every action $n\ge2$ extends from pure monomials to
all nonconstant filters
$\Pi(u)=c\,p(u/\pi)$, where $c\in\R\setminus\{0\}$ and $p$ has
real algebraic coefficients. At integer or half-integer centers no such
filter can cancel even one of the two data at action two. In particular
arbitrarily high degree in this coefficient class cannot reproduce the
deep logarithmic-node filters of Theorem~\ref{ojf:thm:cost}.
\end{corollary}
\begin{proof}
The real and imaginary parts of $p(-a+\iu)$ are nonzero polynomials in
$a$ with real algebraic coefficients: their degrees are respectively
$\deg p$ and $\deg p-1$. A number transcendental over $\Q$ cannot be
a root of a nonzero polynomial over the algebraic numbers. Apply
Lemma~\ref{ojf:lem:log-transcendence} and then the integer or half-integer
rotation. The common nonzero real scale $c$ does not affect vanishing.
\end{proof}

This use of Gelfond--Schneider proves an all-order nonvanishing statement;
a finite table of numerical coefficients would not prove it. Conversely,
the freely chosen real coefficients in Theorem~\ref{ojf:thm:joint} deliberately
permit cancellation. There is no contradiction: a pure monomial in the
order derivative is a much smaller class of filters.

\subsubsection{A normalized sharp jump filter}
For $N\ge1$ define
\begin{equation}\label{ojf:eq:QN}
 Q_N(u)=\prod_{j=1}^N(u+\log j).
\end{equation}
There is a unique real polynomial $\Pi_N$ with
\begin{equation}\label{ojf:eq:PiN}
 \sin(\pi D)\Pi_N=Q_N,\qquad
      (\cos(\pi D)\Pi_N)(0)=1.
\end{equation}
It has degree $N+1$. Existence and uniqueness follow from the
isomorphism modulo constants used in Theorem~\ref{ojf:thm:cost}; the second
condition fixes that constant. One constructive formula is
\[
 \Pi_N=\frac1\pi\frac{\pi D}{\sin(\pi D)}
               \int_0^u Q_N(v)\dd v+c_N,
\]
where $c_N$ enforces the normalization. All operators terminate.
Writing $R_N=\cos(\pi D)\Pi_N$, the first two examples are
\begin{align}
 \Pi_1(u)&=\frac{u^2}{2\pi}+1+\frac\pi2,
 &R_1(u)&=1+\frac{u^2}{2\pi},\label{ojf:eq:Pi1}\\
 \Pi_2(u)&=\frac{u^3}{3\pi}+
       \frac{\log2}{2\pi}u^2+\frac\pi3u+1+\frac{\pi\log2}{2},
 &R_2(u)&=1+\frac{u^3}{3\pi}+
       \frac{\log2}{2\pi}u^2-\frac{2\pi}{3}u.
                                                        \label{ojf:eq:Pi2}
\end{align}
The normalization is intentionally about the \emph{completion}, not
about $\Pi_N(0)$.

\begin{theorem}[Convergent Stokes-neutral curve and its crossover]
\label{ojf:thm:neutral}
Fix an integer $m\ge0$ and $N\ge1$. For sufficiently small $q$, there
is a unique holomorphic function $\epsilon_N(q)$ with
$\epsilon_N(0)=0$ for which
\begin{equation}\label{ojf:eq:neutral-identity}
 J_{m+\epsilon_N(q),\Pi_N}(-\log q)=0
\end{equation}
when $q>0$. For complex $q$, this identity refers to the convergent
$q$-series of the jump. The function has real Taylor coefficients and
\begin{equation}\label{ojf:eq:neutral-leading}
 \epsilon_N(q)=-\frac{\kappa_{m,N}}\pi q^N+O(q^{N+1}),\qquad
 \kappa_{m,N}=\frac1{(N+1)^m}
                 \prod_{j=1}^N\log\frac{N+1}{j}>0.
\end{equation}
The associated completion does \emph{not} disappear:
\begin{equation}\label{ojf:eq:neutral-completion}
 C_{m+\epsilon_N(q),\Pi_N}(-\log q)
                  =(-1)^m q+O(q^2).
\end{equation}
More generally, with $\epsilon=\lambda q^N$,
\begin{equation}\label{ojf:eq:crossover}
 \frac{J_{m+\lambda q^N,\Pi_N}(-\log q)}
                   {-2\iu(-1)^m q^{N+1}}
       =\pi\lambda+\kappa_{m,N}+O(q),
\end{equation}
uniformly for $\lambda$ in any fixed compact subset of $\C$.
At $m=0$ these are analytically continued lateral realizations; no
unregularized order-zero boundary integral is assumed.
\end{theorem}

\begin{proof}
Remove the factor $-2\iu(-1)^m q$ from the jump and define
\begin{align}
 \Psi(\epsilon,q)=\sum_{n\ge1}(-1)^{n-1}n^{-m-\epsilon}q^{n-1}
 \{\sin(\pi\epsilon)R_N(u_n)+
                  \cos(\pi\epsilon)Q_N(u_n)\}.\label{ojf:eq:Psi}
\end{align}
The series is holomorphic for $|q|<1$ and $\epsilon$ in a neighborhood
of zero: on compact sets its coefficients have at most a fixed power of
$n$ times a fixed power of $\log n$. Its $n=1$ term is precisely
$\sin(\pi\epsilon)$, since $Q_N(0)=0$ and $R_N(0)=1$.
Consequently
\[
 \Psi(0,0)=0,\qquad \partial_\epsilon\Psi(0,0)=\pi.
\]
The holomorphic implicit-function theorem gives the unique curve.
Conjugation and uniqueness imply real Taylor coefficients.

At $\epsilon=0$, the first $N$ coefficients vanish. At $n=N+1$,
\[
 (-1)^N(N+1)^{-m}Q_N(-\log(N+1))=\kappa_{m,N}.
\]
Thus $\Psi(0,q)=\kappa_{m,N}q^N+O(q^{N+1})$. Substitution in the
implicit equation, whose $\epsilon$ derivative is $\pi+O(q)$, gives
\eqref{ojf:eq:neutral-leading}. Taylor expansion with
$\epsilon=\lambda q^N$ proves \eqref{ojf:eq:crossover}; higher sine
terms and mixed terms are of sufficiently high order, and holomorphy
gives uniformity on compact $\lambda$ sets.

The corresponding completion has normalized coefficient
\[
 \cos(\pi\epsilon)R_N(u_n)-\sin(\pi\epsilon)Q_N(u_n).
\]
Its first term is $(-1)^m q\cos(\pi\epsilon)$. The rest is $O(q^2)$
uniformly for small $\epsilon$, and the implicit curve tends to zero.
This proves \eqref{ojf:eq:neutral-completion}.
\end{proof}

\begin{remark}[What the curve cancels]
For a \emph{fixed} order, the sharp filters remove individual initial
action coefficients. On the moving-order curve, coefficients at different
actions compensate after $\epsilon$ is substituted as a function of $q$.
The result is an exact identity of analytic functions, not a statement
that every fixed-order Borel singularity has disappeared. It also gives
a precise example in which zero total lateral jump does not identify the
physical function with the common lateral realization.
\end{remark}

\subsection{Transport along the exact neutral curve}
\label{ojf:sec:neutral-inverse}

The neutral curve cancels the entire lateral difference by allowing the
order to move with the exponential parameter. Inversion must retain this
dependence. The following result answers the inverse-transport question
in \cite{ojf:report} for every fixed filter size and nonnegative integer
center; it also supplies an exact convergent coefficient formula.

\begin{theorem}[A convergent neutral inverse chart and its physical displacement]
\label{ojf:thm:neutral-inverse}
Fix $N\ge1$ and $m\in\mathbb Z_{\ge0}$, and use the normalized filter
$\Pi_N$ and analytic function $\varepsilon_N$ of
Theorem~\ref{ojf:thm:neutral}, with this order center. Put
\[
 \sigma(x)=m+\varepsilon_N(e^{-x}),\qquad
 P(x)=\left.\Pi_N(\partial_s)F_s(x)\right|_{s=\sigma(x)},\qquad
 L^\pm(x)=\left.\Pi_N(\partial_s)L_s^\pm(x)\right|_{s=\sigma(x)}.
\]
For all sufficiently large real $x$, the two lateral functions coincide:
$L^+(x)=L^-(x)=L(x)$. The functions $P$ and $L$ are real analytic,
strictly increasing and unbounded on a sufficiently late half-line.
Let $X$ denote the inverse of $P$ there. Write
\[
 C_w=L(w)-P(w),\qquad
 H_w(t)=\begin{cases}(P(w+t)-P(w))/t,&t\ne0,\\P'(w),&t=0.\end{cases}
\]
For any fixed $0<\rho<1$, and all sufficiently large $w$, the series
\begin{equation}\label{ojf:eq:neutral-inverse-chart}
 X(P(w)+zC_w)-w
   =\sum_{k\ge1}\frac{(zC_w)^k}{k}
                 [t^{k-1}]H_w(t)^{-k}
\end{equation}
converges for $|z|<\rho P'(w)/(2|C_w|)$, where the left side denotes
the local holomorphic inverse extending the real one. This disk contains
$z=1$ for large $w$. In particular, the common lateral inverse at target
$L(w)$ equals $w$, whereas the physical inverse satisfies
\begin{equation}\label{ojf:eq:neutral-inverse-shift}
 X(L(w))-w\sim\frac{(-1)^m e^{-w}}{P'(w)}
 \sim
 \begin{cases}
 \displaystyle\frac{\pi w e^{-w}}{(\log w)^N},&m=0,\\[6pt]
 \displaystyle\frac{(-1)^m\pi(N+1)(m-1)!\,w^{1-m}e^{-w}}
                         {(\log w)^{N+1}},&m\ge1.
 \end{cases}
\end{equation}
The first two exact terms of \eqref{ojf:eq:neutral-inverse-chart} are
$zC_w/P'(w)-z^2P''(w)C_w^2/(2P'(w)^3)$.
\end{theorem}

\begin{proof}
The neutral-curve identity proves equality of the laterals and gives
\begin{equation}\label{ojf:eq:neutral-C-asymptotic}
 C(x):=L(x)-P(x)=(-1)^m e^{-x}+O(e^{-2x}).
\end{equation}
The underlying normally convergent power series in $q$ and the analytic
implicit function $\varepsilon_N(q)$ also give the differentiated
estimate and its holomorphic version near the positive ray.

Set $d=N+1$. The leading coefficient of $\Pi_N$ is $1/(\pi d)$,
since $\sin(\pi D)$ lowers degree by one with leading multiplier
$\pi d$ and its image is the monic $Q_N$. Uniformly for $s$ in a small
fixed disk about $m$, the Hurwitz realization
\eqref{ojf:eq:hurwitz-completion} gives the finite Sommerfeld estimate
\[
 F_s(x)=\frac{x^s}{\Gamma(s+1)}+O(|x|^{m+\eta-2}),
 \qquad |s-m|\le\eta<\tfrac12,
\]
on a smaller sector about the positive ray. Here the exponential
completion has been absorbed in the remainder. To obtain the estimate
and any of its finitely many derivatives, take a slightly larger compact
order disk in the sectorial Hurwitz expansion
\cite[Eq.~25.11.43]{intdist:DLMFHurwitz}, then apply Cauchy's estimates
in the order and in $x$. The odd $x^{s-1}$ term cancels after the shift
$1/2$ in the Hurwitz argument; equivalently the first correction in
\eqref{ojf:eq:sommerfeld} has power $x^{s-2}$. This finite estimate,
not termwise differentiation of a divergent infinite series, suffices.

At the fixed center $m$, applying $\Pi_N(\partial_s)$ to the leading
term gives $x^m/m!$ times a polynomial in $\log x$ of degree $d$,
with leading coefficient $1/(\pi d)$. The chain rule for the moving
center is
\begin{equation}\label{ojf:eq:moving-chain-rule}
 P'(x)=\left.\Pi_N(\partial_s)\partial_xF_s(x)\right|_{s=\sigma(x)}
       +\sigma'(x)\left.\Pi_N(\partial_s)\partial_sF_s(x)
                                                       \right|_{s=\sigma(x)},
 \qquad \sigma'(x)=-e^{-x}\varepsilon_N'(e^{-x}).
\end{equation}
Both $\sigma-m$ and $\sigma'$ are $O(e^{-Nx})$; compact-order
Cauchy estimates bound the terms they introduce by an exponential
times a fixed power of $x$. They are therefore negligible relative to
the leading derivative. Thus
\begin{equation}\label{ojf:eq:neutral-derivative}
 P'(x)\sim
 \begin{cases}
 (\log x)^N/(\pi x),&m=0,\\
 x^{m-1}(\log x)^{N+1}/\bigl(\pi(N+1)(m-1)!\bigr),&m\ge1.
 \end{cases}
\end{equation}
The same sectorial estimates hold for derivatives and imply
$P'(w+t)/P'(w)\to1$ uniformly on every fixed complex disk in $t$.
The leading expression for $P$ tends to $+\infty$ also when $m=0$.
Together with \eqref{ojf:eq:neutral-C-asymptotic}, this proves eventual
monotonicity, unboundedness and existence of both real inverses.

Uniformly for $|t|\le\rho$, integration of the derivative estimate
gives $H_w(t)/P'(w)\to1$. For large $w$ its modulus is greater than
$1/2$. On $|t|=\rho$ the perturbation $zC_w$ is therefore smaller
than $P(w+t)-P(w)$ whenever $|z|<\rho P'(w)/(2|C_w|)$.
Rouch\'e's theorem gives exactly one zero, counted with multiplicity,
of $P(w+t)-P(w)-zC_w$ in the disk. It is simple and varies
holomorphically with $z$. Lagrange inversion yields
\eqref{ojf:eq:neutral-inverse-chart} throughout this disk.

Finally $|C_w|/P'(w)\to0$, so its radius in $z$ tends to infinity.
At $z=1$ the solution $t_w$ obeys
$t_wH_w(t_w)=C_w$, hence
$t_w\sim C_w/P'(w)$. Substitution of
\eqref{ojf:eq:neutral-C-asymptotic} and
\eqref{ojf:eq:neutral-derivative} proves the displayed displacement.
Expanding $H_w(t)=P'(w)+P''(w)t/2+\cdots$ gives the second coefficient.
\end{proof}

Thus an exactly vanishing lateral difference can coexist with a nonzero
first exponential action in the physical inverse. The sign of that
displacement alternates with $m$. The theorem concerns fixed $m,N$ and
an eventual inverse chart; it does not assert a growing-filter limit or
an exact global complex inverse radius.


\subsection{All order-zero jets in Stieltjes and integral coordinates}
\label{ojf:sec:stieltjes}

\subsubsection{A complete finite formula in generalized Stieltjes constants}
For $\Real a>0$ use the Hurwitz--Stieltjes convention \cite{intdist:DLMFHurwitz}
\begin{equation}\label{ojf:eq:stieltjes}
 \zeta(1+z,a)=\frac1z+
       \sum_{n\ge0}\frac{(-1)^n\gamma_n(a)}{n!}z^n.
\end{equation}
In particular $\gamma_0(a)=-\psi(a)$.
Define the complete exponential Bell polynomials by
\begin{equation}\label{ojf:eq:bell}
 \exp\left(\sum_{j\ge1}b_j\frac{z^j}{j!}\right)
       =\sum_{r\ge0}\mathcal{B}_r(b_1,\ldots,b_r)\frac{z^r}{r!},
 \qquad\mathcal{B}_0=1.
\end{equation}
The following theorem specializes classical Jonqui\`ere inversion to a
resonance excluded from the positive-integer Gamma-kernel statement of
\cite{qfb:reports}.

\begin{theorem}[Stieltjes coordinates for every zero-order lateral jet]
\label{ojf:thm:stieltjes}
Let $x>0$, $a_\pm=\frac12\pm\frac{\iu x}{2\pi}$, and
\begin{align}
 d_1^\pm&=\gamma+\log(2\pi)\mp\frac{\iu\pi}{2},\label{ojf:eq:dj}\\
 d_j^\pm&=(-1)^{j+1}(j-1)!\zeta(j),\qquad j\ge2.\notag
\end{align}
For $r\ge1$ put
\begin{equation}\label{ojf:eq:stieltjes-jet}
 Z_r^\pm(x)=\mathcal{B}_r(d^\pm)
       -r\sum_{n=0}^{r-1}\binom{r-1}{n}
                    \gamma_n(a_\pm)\mathcal{B}_{r-1-n}(d^\pm),
 \qquad Z_0^\pm=1.
\end{equation}
Then $Z_r^\pm=\partial_s^rL_s^\pm|_{s=0}$ and
\begin{equation}\label{ojf:eq:physical-zero-jets}
 \left.\partial_s^rF_s(x)\right|_{s=0}
  =Z_r^\pm(x)-\sum_{n\ge1}(-1)^{n-1}
                     (-\log n\mp\iu\pi)^re^{-nx}.
\end{equation}
The formulas hold for every finite $r$ and locally uniformly for $x>0$.
\end{theorem}

\begin{proof}
The reciprocal-Gamma expansion \cite[\S5.7]{ojf:DLMFGamma} gives
\[
 \frac1{\Gamma(s)}=s\exp\left(\gamma s+
           \sum_{j\ge2}\frac{(-1)^{j+1}\zeta(j)}j s^j\right),
 \qquad |s|<1.
\]
In \eqref{ojf:eq:stieltjes} substitute $z=-s$. Combining it with
\eqref{ojf:eq:hurwitz-completion} yields the exact local generating
identity
\begin{equation}\label{ojf:eq:stieltjes-generating}
 L_s^\pm(x)=
 \exp\left(\sum_{j\ge1}d_j^\pm\frac{s^j}{j!}\right)
       \left(1-\sum_{n\ge0}\gamma_n(a_\pm)\frac{s^{n+1}}{n!}\right).
\end{equation}
In particular the removable value at zero is $1$, not zero. Multiplying
Taylor series and taking $r!$ times the coefficient gives
\eqref{ojf:eq:stieltjes-jet}. Applying the exact completion formula at
zero gives \eqref{ojf:eq:physical-zero-jets}. The local Taylor identities
and the normally convergent exponential series justify every finite
coefficient extraction and the claimed local uniformity.
\end{proof}

For example,
\[
 Z_1^\pm=d_1^\pm-\gamma_0(a_\pm),\qquad
 Z_2^\pm=(d_1^\pm)^2-\zeta(2)
                   -2d_1^\pm\gamma_0(a_\pm)-2\gamma_1(a_\pm).
\]
The order-zero second physical derivative is therefore an explicit
combination of $\gamma_1$ at a conjugate pair of points, digamma values,
and a convergent log-weighted exponential series. There is no finite-part
ambiguity in these identities: the regularization is fixed by the
holomorphic product in \eqref{ojf:eq:hurwitz-completion}.

\subsubsection{Independent endpoint-subtracted integral coordinates}
A second exact representation is useful both for proof and for independent
numerical evaluation. It avoids differentiating a numerically delicate
polylogarithm continuation at an integer order.

\begin{theorem}[Positive regularization and all zero-order moments]
\label{ojf:thm:regularization}
For any real $x$ let $q=e^{-x}$ and $f_0=(1+q)^{-1}$. Define
\begin{equation}\label{ojf:eq:moments}
 M_j(x)=\frac1{1+q}\int_0^\infty
       \frac{(1-e^{-t})(\log t)^j}{t(1+e^{t-x})}\dd t,
       \qquad j\ge0.
\end{equation}
All these integrals are absolutely convergent. With
$g_1=\gamma$ and $g_j=(-1)^{j+1}(j-1)!\zeta(j)$ for $j\ge2$,
\begin{equation}\label{ojf:eq:moment-jets}
 F_0(x)=f_0,\qquad
 \left.\partial_s^rF_s(x)\right|_{s=0}
   =r\sum_{j=0}^{r-1}\binom{r-1}{j}\mathcal{B}_{r-1-j}(g)M_j(x)
       \quad(r\ge1).
\end{equation}
In particular $H(x):=\partial_sF_s(x)|_{s=0}=M_0(x)>0$.
\end{theorem}

\begin{proof}
Let $f_x(t)=(1+e^{t-x})^{-1}$. The subtraction has the exact
positive form
\begin{equation}\label{ojf:eq:positive-subtraction}
 f_x(t)-f_0e^{-t}
  =\frac{1-e^{-t}}{(1+e^{-x})(1+e^{t-x})}.
\end{equation}
For $\Real s>0$, subtract the elementary Gamma integral to get
\begin{equation}\label{ojf:eq:regularized-F}
 F_s(x)=f_0+\frac1{\Gamma(s)}\int_0^\infty
                          t^{s-1}(f_x(t)-f_0e^{-t})\dd t.
\end{equation}
At zero the subtracted numerator is $O(t)$, so the remaining integral
is holomorphic for $\Real s>-1$. At infinity it has exponential decay
on compact real $x$ sets. Multiplication by $1/\Gamma(s)$ continues
\eqref{ojf:eq:regularized-F} to a neighborhood of zero.
Its $j$th integral derivative at zero is $M_j(x)$; the majorants
$t^{-1+\epsilon}(1+|\log t|^j)$ after subtraction near zero, with
$\epsilon>0$, and exponential majorants at infinity justify this.
The reciprocal-Gamma Bell expansion used above now gives
\eqref{ojf:eq:moment-jets}. Positivity of the integrand for $j=0$ proves
the final assertion.
\end{proof}

\subsection{The order-zero digamma core and its exact Borel kernel}
\label{ojf:sec:borelzero}

\subsubsection{An exact physical identity}
Put
\begin{align}
 A(x)&=\gamma+\log(2\pi)+
              \Real\psi\left(\tfrac12+\frac{\iu x}{2\pi}\right),
                   &&x>0,\label{ojf:eq:A-real}\\
 B(z)&=\sum_{n\ge2}(-1)^n(\log n)z^n
          =(\log2)z^2-(\log3)z^3+(\log4)z^4-\cdots,
                   &&|z|<1.\label{ojf:eq:B}
\end{align}
The letter $B$ here denotes a convergent analytic function, not the
Borel transform $\mathcal{B}$.

\begin{proposition}[Digamma core and both lateral first jets]
\label{ojf:prop:digamma}
For $x>0$,
\begin{align}
 H(x)&=A(x)-B(e^{-x}),\label{ojf:eq:H-A}\\
 Z_1^\pm(x)&=A(x)\mp\iu\pi\frac{e^{-x}}{1+e^{-x}}.
                                                        \label{ojf:eq:Z1}
\end{align}
Thus the first zero-order jet has a jump beginning at action one,
while its physical correction to the arithmetic average begins at
action two.
\end{proposition}
\begin{proof}
By Theorem~\ref{ojf:thm:stieltjes},
$Z_1^+=\gamma+\log(2\pi)-\iu\pi/2+\psi(a_+)$.
The digamma reflection formula \cite[\S5.5]{ojf:DLMFGamma} gives
\[
 \psi(1-z)-\psi(z)=\pi\cot(\pi z),\qquad
 \Imag\psi\left(\tfrac12+\frac{\iu x}{2\pi}\right)
                =\frac\pi2\tanh(x/2).
\]
Hence the imaginary part of $Z_1^+$ is $-\pi/(e^x+1)$, and
its real part is $A$. The lower formula follows by conjugation.
In \eqref{ojf:eq:physical-zero-jets} with $r=1$, the imaginary tail
cancels exactly, while the real tail is $B(e^{-x})$.
\end{proof}

\begin{corollary}[Symmetry and the global positive core branch]
\label{ojf:cor:core-global}
The even function $A$ satisfies
\[
 H(x)+H(-x)=A(x),\qquad A(0)=\log(\pi/2),
 \qquad H(0)=\tfrac12\log(\pi/2).
\]
It is strictly increasing for $x>0$ and maps that half-line onto
$(\log(\pi/2),\infty)$.
\end{corollary}
\begin{proof}
For $x>0$, $H(-x)=B(e^{-x})$, giving the symmetry from
\eqref{ojf:eq:H-A}; continuity and evenness give the remaining real
points. The classical value $\psi(1/2)=-\gamma-2\log2$ gives $A(0)$.
For $v=x/(2\pi)>0$, the convergent trigamma series gives
\[
 A'(x)=-\frac1{2\pi}\Imag\sum_{k\ge0}\frac1{(k+1/2+\iu v)^2}>0,
\]
because each summand has strictly negative imaginary part. The range
then follows from the logarithmic asymptotic proved below. This also
shows that the exact positive core inverse used later is unique wherever
its target exceeds $\log(\pi/2)$.
\end{proof}

The single meromorphic expression, holomorphic near large real points, is
\begin{equation}\label{ojf:eq:A-holomorphic}
 A(z)=\gamma+\log(2\pi)+\frac12
    \left\{\psi\left(\tfrac12+\frac{\iu z}{2\pi}\right)
          +\psi\left(\tfrac12-\frac{\iu z}{2\pi}\right)\right\}.
\end{equation}
It is meromorphic in $z$, with possible poles on the imaginary axis,
and agrees with \eqref{ojf:eq:A-real} on the positive real axis. In
particular it is holomorphic on every disk of a fixed radius around a
sufficiently large positive real $w$. We use this expression, not the
nonholomorphic operation $\Real$, in complex inverse equations.

\subsubsection{A meromorphic minor at the previously excluded order}

\begin{theorem}[Ordinary Borel reconstruction at order zero]
\label{ojf:thm:borelzero}
The formal algebraic tail of $A(x)-(\log x+\gamma)$ is
\begin{equation}\label{ojf:eq:A-asymptotic}
 A(x)\sim\log x+\gamma-
          2\sum_{k\ge1}\eta(2k)(2k-1)!x^{-2k}.
\end{equation}
Its ordinary Borel minor is the meromorphic function
\begin{equation}\label{ojf:eq:b0}
 b_0^{\rm jet}(\xi)=\frac1\xi-\frac\pi{\sin\pi\xi},
\end{equation}
with a removable zero at the origin and residue $(-1)^{n-1}$ at
$\xi=n\ge1$. For $x>0$,
\begin{align}
 \log x+\gamma+\int_\pm e^{-x\xi}b_0^{\rm jet}(\xi)\dd\xi
     &=Z_1^\pm(x),\label{ojf:eq:b0-lateral}\\
 A(x)&=\log x+\gamma+
          \operatorname{PV}\int_0^\infty e^{-x\xi}b_0^{\rm jet}(\xi)\dd\xi.
                                                        \label{ojf:eq:b0-PV}
\end{align}
The principal value is the arithmetic average of the lateral integrals,
or the sum of symmetric local principal values at the positive integers.
\end{theorem}

\begin{proof}
The classical generator
$\pi\xi/\sin\pi\xi=1+2\sum_{k\ge1}\eta(2k)\xi^{2k}$
shows that
\[
 b_0^{\rm jet}(\xi)=-2\sum_{k\ge1}\eta(2k)\xi^{2k-1}.
\]
The ordinary Borel convention gives exactly the coefficients in
\eqref{ojf:eq:A-asymptotic}. The residue statement follows from
$(\sin\pi\xi)'|_{\xi=n}=\pi(-1)^n$.

For an independent proof of the exact reconstruction, initially take
$\Real x>0$ and $\Imag x<0$. Rotate an upper tilted Laplace path to
$\xi=\iu t$. There are no poles inside the rotation sector. On its
large arcs, the cosecant has exponential decay away from the real
axis and the $1/\xi$ term is suppressed by $e^{-x\xi}$; the negative
imaginary part of $x$ ensures decay even at the final imaginary ray.
The removable origin gives no small-circle contribution. Thus
\begin{align*}
 \int_+e^{-x\xi}b_0^{\rm jet}(\xi)\dd\xi
 &=\int_0^\infty e^{-\iu xt}
                   \left(\frac1t-\frac\pi{\sinh\pi t}\right)\dd t\\
 &=\int_0^\infty
       \left\{\frac{e^{-\iu xu/(2\pi)}}u
          -\frac{e^{-(1/2+\iu x/(2\pi))u}}{1-e^{-u}}\right\}\dd u.
\end{align*}
The standard digamma integral, combined with the Frullani difference,
is \cite[\S5.9]{ojf:DLMFDigamma}
\[
 \psi(a)=\int_0^\infty
          \left(\frac{e^{-u}}u-\frac{e^{-au}}{1-e^{-u}}\right)\dd u,
 \quad \Real a>0,
\]
where cancellation at zero is taken inside the integrand. Consequently
the preceding integral is
\begin{equation}\label{ojf:eq:b0-evaluation}
 \psi\left(\tfrac12+\frac{\iu x}{2\pi}\right)
                 -\log x+\log(2\pi)-\frac{\iu\pi}{2}.
\end{equation}
Both sides continue to a neighborhood of the positive real axis. A fixed
upper ray of angle less than $\pi/2$ has an absolutely convergent
Laplace integral there, so \eqref{ojf:eq:b0-evaluation} proves the upper
formula for real positive $x$. Conjugation proves the lower formula.
Its value is precisely \eqref{ojf:eq:Z1}.

On the real boundary the only singularities are simple poles. Upper
minus lower indentation contributes minus $2\pi\iu$ times their
residues, and the weighted residue sum is absolutely convergent:
\[
 -2\pi\iu\sum_{n\ge1}(-1)^{n-1}e^{-nx}
          =-2\pi\iu\frac{e^{-x}}{1+e^{-x}}.
\]
At infinity one may take boundaries midway between the integer poles;
exponential decay controls the remaining contour and principal-value
tails. Averaging proves \eqref{ojf:eq:b0-PV}.
Finally, a finite Taylor expansion of the germ on a tilted ray, with a
Laplace remainder, gives \eqref{ojf:eq:A-asymptotic} to each finite order.
\end{proof}

The kernel \eqref{ojf:eq:b0} is the natural $m=0$ extension of the
subtracted Gamma kernel in \cite[label
\texttt{thm:Gamma-kernel}]{qfb:reports}. The extension has been proved
here with a separate contour calculation, rather than by asserting
that all hypotheses of the $m\ge1$ theorem persist unchanged.

\begin{corollary}[Derivative scales for the exact core]
\label{ojf:cor:core-derivatives}
As $w\to+\infty$,
\begin{align}
 A(w)&=\log w+\gamma-\frac{\pi^2}{6w^2}+O(w^{-4}),\label{ojf:eq:A-two}\\
 A'(w)&=\frac1w+O(w^{-3}),\qquad
 A''(w)=-\frac1{w^2}+O(w^{-4}).\label{ojf:eq:A-derivatives}
\end{align}
The estimates are holomorphically uniform in a fixed closed subsector
about the positive real axis, away from the origin. In particular,
$A'(w)>0$ for all sufficiently large real $w$.
\end{corollary}
\begin{proof}
Use the first term of \eqref{ojf:eq:A-asymptotic}. Uniformity follows
from the tilted Laplace proof, or the standard sectorial digamma
expansion \cite[Eq.~5.11.2]{ojf:DLMFDigamma}, applied to the conjugate-pair
formula \eqref{ojf:eq:A-holomorphic}. Cauchy estimates on disks of radius
a fixed positive multiple of $w$, contained in such a subsector, justify
differentiation of each finite asymptotic remainder and give
\eqref{ojf:eq:A-derivatives}. This uses holomorphic remainder control,
not mere differentiation of a real-axis big-$O$ statement.
\end{proof}

\subsection{A global inverse and convergent charts at both endpoints}
\label{ojf:sec:inverse}

\subsubsection{The real inverse exists globally}

\begin{theorem}[Global order-zero inverse]\label{ojf:thm:global}
The function $H(x)=\partial_sF_s(x)|_{s=0}$ is a real-analytic strictly
increasing bijection from $\R$ to $(0,\infty)$. Its inverse $X(y)$ is
real analytic on $(0,\infty)$. Its two endpoint scales are
\begin{equation}\label{ojf:eq:H-endpoints}
 H(x)\sim(\log2)e^{2x}\quad(x\to-\infty),\qquad
 H(x)\sim\log x\quad(x\to+\infty).
\end{equation}
\end{theorem}
\begin{proof}
Use the positive integral $H=M_0$ from
Theorem~\ref{ojf:thm:regularization}. Both factors
$(1+e^{-x})^{-1}$ and $(1+e^{t-x})^{-1}$ increase strictly with
$x$ for every $t>0$. Differentiation on compact $x$ intervals is
justified by the same endpoint and exponential majorants, and the
integrand derivative is strictly positive. Hence $H'(x)>0$.
Local holomorphic continuation of the integral in $x$ in a strip of
sufficiently small width gives real analyticity.

For $x<0$, the original small-argument polylogarithm series and normal
order differentiation give $H(x)=B(e^x)$, so
$H(x)\sim(\log2)e^{2x}\to0$. For large positive $x$,
\eqref{ojf:eq:H-A} and \eqref{ojf:eq:A-two} give
$H(x)=\log x+\gamma+O(x^{-2})\to\infty$.
Continuity and strict increase prove the range assertion, and the
analytic inverse-function theorem applies everywhere because $H'>0$.
\end{proof}

\subsubsection{The small-target logarithmic--square-root chart}
Write $a=\log2$ and
\begin{equation}\label{ojf:eq:small-b}
 b(z)=\frac{B(z)}{az^2}
       =1-\frac{\log3}{a}z+\frac{\log4}{a}z^2-\cdots,
 \qquad b(0)=1.
\end{equation}
The square-root parameter for a small target is
$t=(y/a)^{1/2}>0$.

\begin{theorem}[All coefficients and an explicit convergent endpoint disk]
\label{ojf:thm:small-inverse}
For $0<t<\sqrt3/40$, the global inverse has the convergent expansion
\begin{equation}\label{ojf:eq:small-inverse}
 X(at^2)=\log t+\sum_{N\ge1}c_Nt^N,
 \qquad c_N=\frac1N\coef{z}{N}b(z)^{-N/2}.
\end{equation}
The analytic correction $X(at^2)-\log t$ extends holomorphically to
$|t|<\sqrt3/40$. In particular,
\begin{equation}\label{ojf:eq:small-first}
 X(y)=\frac12\log\frac ya+
     \frac{\log3}{2a}\left(\frac ya\right)^{1/2}
     +\left(\frac{(\log3)^2}{2a^2}-1\right)\frac ya
     +O(y^{3/2}).
\end{equation}
The radius given here is a sufficient certificate, not an exact radius.
\end{theorem}

\begin{proof}
Let $r=1/20$. Since $\log n\le n$ for $n\ge1$,
\[
 \sup_{|z|\le r}|b(z)-1|
 \le\frac1a\sum_{n\ge3}n r^{n-2}
 =\frac{r(3-2r)}{a(1-r)^2}
 =\frac{58}{361a}<\frac14.
\]
For the strict last inequality, $\log2>2/3$ follows, for example, from
$\log2=2\sum_{k\ge0}(1/3)^{2k+1}/(2k+1)$; then
$58/(361a)<87/361<1/4$.
There is a holomorphic square root of $b$ on the disk, equal to one
at zero, and $|\sqrt{b(z)}|>\sqrt3/2$ there. The function
$f(z)=z\sqrt{b(z)}$ has exactly the simple zero $0$, and on $|z|=r$
its modulus exceeds $\sqrt3/40$.

For $|t|<\sqrt3/40$, Rouch\'e's theorem gives exactly one root,
counting multiplicity, of $f(z)-t$ inside $|z|<r$. That root is simple,
so it defines a holomorphic function $z(t)$ with $z(0)=0$ and
$z'(0)=1$. Uniqueness glues the local roots throughout the disk.
The quotient $z(t)/t$ extends holomorphically and has no zeros;
its logarithm with value zero at $t=0$ therefore exists on the disk.
For positive small $t$, the root is real and positive and
$H(\log z(t))=B(z(t))=at^2$. The global uniqueness theorem identifies
$\log z(t)$ with $X(at^2)$.

For the coefficient formula, set $\Phi(z)=b(z)^{-1/2}$. The equation
is $z=t\Phi(z)$ and $\log(z/t)=\log\Phi(z)$.
Lagrange--B\"urmann inversion gives, for $N\ge1$,
\begin{align*}
 \coef{t}{N}\log(z(t)/t)
 &=\frac1N\coef{z}{N-1}(\log\Phi)'\Phi^N\\
 &=\frac1{N^2}\coef{z}{N-1}(\Phi^N)'
  =\frac1N\coef{z}{N}\Phi^N.
\end{align*}
This proves \eqref{ojf:eq:small-inverse}. Expanding the first two
coefficients and using $\log4=2\log2$ proves
\eqref{ojf:eq:small-first}.
\end{proof}

\subsubsection{The large-target exact-core chart}
For large $w>0$, take the target $y=A(w)$, so $w$ is the exact
inverse of the algebraic digamma core rather than an approximation to
that inverse. Put
\begin{equation}\label{ojf:eq:large-notations}
 q=e^{-w},\qquad a_w=A'(w),\qquad b_w=A''(w),\qquad
 c_+=-\iu\pi,\quad c_-=\iu\pi.
\end{equation}
Let $G_\pm(y)$ denote the inverse branches near $w$ of the lateral
functions $Z_1^\pm(x)$ from \eqref{ojf:eq:Z1}.

\begin{theorem}[Convergent inverse sectors and the exact second sector]
\label{ojf:thm:large-inverse}
There are constants $c>0$ and $w_0>0$ such that for each $w>w_0$ the
three inverse displacements are holomorphic for $|q|<c/w$ when
$q$ is treated as an independent variable. At the physical substitution
$q=e^{-w}$, their convergent expansions give
\begin{align}
 X(A(w))&=w+\frac{\log2}{a_w}q^2+O(w^3q^3),
                                            \label{ojf:eq:physical-inverse}\\
 G_\pm(A(w))&=w-\frac{c_\pm}{a_w}q+
 \left(-\frac{b_wc_\pm^2}{2a_w^3}
       -\frac{c_\pm^2}{a_w^2}+\frac{c_\pm}{a_w}\right)q^2
                   +O(w^3q^3).              \label{ojf:eq:lateral-inverse}
\end{align}
The remainder constants can be chosen independently of large $w$.
The physical branch in \eqref{ojf:eq:physical-inverse} is the global
inverse of Theorem~\ref{ojf:thm:global}.
\end{theorem}

\begin{proof}
For independent $q$, the defining equations for the displacements are
\begin{align}
 A(w+\delta)-A(w)-B(qe^{-\delta})&=0,
                                             \label{ojf:eq:inv-physical-eq}\\
 A(w+\delta)-A(w)+c_\pm
       \frac{qe^{-\delta}}{1+qe^{-\delta}}&=0.
                                             \label{ojf:eq:inv-lateral-eq}
\end{align}
Fix a small radius $\rho>0$, for example $\rho=1/4$. The uniform
estimates in Corollary~\ref{ojf:cor:core-derivatives} imply
\[
 A(w+\delta)-A(w)=\frac\delta w(1+O(w^{-1}))
                    \qquad (|\delta|\le\rho).
\]
For sufficiently large $w$, the core difference has exactly one zero
in the disk and has boundary modulus at least $\rho/(2w)$; this follows
by comparison with $\delta/w$.

On the same disk, the lateral perturbation is bounded by
$\pi |q|e^\rho/(1-|q|e^\rho)$.
Choose $c>0$ small enough that this is less than $\rho/(2w)$ whenever
$|q|\le c/w$, for all sufficiently large $w$.
The physical perturbation $B(qe^{-\delta})$ is $O(|q|^2)$ uniformly
on the disk and satisfies the same strict inequality, after increasing
$w_0$ if necessary. Rouch\'e gives exactly one simple zero in each case.
The implicit-function theorem and uniqueness give holomorphic roots on
the parameter disk, bounded by $\rho$.
Cauchy's estimate bounds their degree-$n$ coefficients by
$\rho(w/c)^n$. Hence on $|q|\le c/(2w)$ the remainder after degree two
is at most
\[
 \rho\frac{(w|q|/c)^3}{1-w|q|/c}=O(w^3|q|^3).
\]
The physical value $q=e^{-w}$ lies in this region for all large $w$.
This proves actual sector convergence and a uniform remainder; it is
not only a formal reversion computation.

For the coefficients, write $\delta=d_1q+d_2q^2+O(q^3)$ and use
\[
 A(w+\delta)-A(w)=a_w\delta+\tfrac12b_w\delta^2+O(\delta^3),
 \qquad
 \frac{qe^{-\delta}}{1+qe^{-\delta}}
       =q-(d_1+1)q^2+O(q^3).
\]
The first coefficient in the lateral equation gives $d_1=-c/a_w$.
The second gives
$a_wd_2+\tfrac12b_wd_1^2-cd_1-c=0$, which is precisely
\eqref{ojf:eq:lateral-inverse}. In the physical equation,
$B(qe^{-\delta})=(\log2)q^2+O(q^3)$, so $d_1=0$ and
$d_2=\log2/a_w$.
For the physical real $q$, conjugation and uniqueness make the physical
root real. Its defining equation is $H(w+\delta)=A(w)$, and the global
monotonicity of $H$ identifies it with $X(A(w))$.
\end{proof}

\subsubsection{A larger inverse error than the physical completion}

\begin{corollary}[Digamma-core arithmetic-averaging defect]
\label{ojf:cor:defect}
With the preceding exact-core normalization,
\begin{align}
 \frac{G_+(A(w))+G_-(A(w))}{2}-X(A(w))
 &=\left(\frac{\pi^2}{a_w^2}
       +\frac{\pi^2b_w}{2a_w^3}-\frac{\log2}{a_w}\right)e^{-2w}
             +O(w^3e^{-3w})                         \label{ojf:eq:defect-exact}\\
 &\sim\pi^2w^2e^{-2w}.                              \label{ojf:eq:defect-leading}
\end{align}
By contrast,
\begin{equation}\label{ojf:eq:physical-leading}
 X(A(w))-w\sim(\log2)we^{-2w},
 \qquad
 G_+(A(w))-G_-(A(w))\sim2\pi\iu we^{-w}.
\end{equation}
In particular the inverse-average error divided by the physical
core-inverse correction tends to infinity, with equivalent
$\pi^2w/\log2$.
\end{corollary}
\begin{proof}
Average \eqref{ojf:eq:lateral-inverse}, using
$c_++c_-=0$ and $c_\pm^2=-\pi^2$, then subtract
\eqref{ojf:eq:physical-inverse}. This gives
\eqref{ojf:eq:defect-exact}. Since
$a_w=w^{-1}(1+O(w^{-2}))$ and
$b_w=-w^{-2}(1+O(w^{-2}))$, the coefficient of $e^{-2w}$ in that
formula is $\pi^2w^2+O(w)$.
The uniform remainder is smaller by a factor $O(we^{-w})$.
The same estimates give the first equivalent in
\eqref{ojf:eq:physical-leading}. Subtracting the two lateral
expansions gives its second equivalent. Division proves the ratio.
\end{proof}

This is a model-specific quantitative consequence of the general
nonlinear averaging mechanism already proved in \cite{ojf:transport}
and restated in \cite{qfb:reports}. What is added here is the explicit
core, the exact second-sector coefficient, a uniform convergent chart,
and the separation by an unbounded factor from the physical correction.
It is \emph{not} a counterexample to a correctly formulated nonlinear
median summation. Nor does it allege that either cited repository article
confused arithmetic averaging with such a summation.

\paragraph{In the target variable.}
Let $v=e^{y-\gamma}$. Inverting \eqref{ojf:eq:A-two} gives
$w=v+\pi^2/(6v)+O(v^{-3})$. Therefore the physical inverse correction
and the inverse-average error have the respective equivalents
\begin{equation}\label{ojf:eq:double-exponential}
 (\log2)e^{y-\gamma}\exp(-2e^{y-\gamma}),\qquad
 \pi^2e^{2(y-\gamma)}\exp(-2e^{y-\gamma}).
\end{equation}
These are doubly exponential scales in $y$. Using the exact $w$ for
finite computations remains preferable: an algebraic approximation to
the core must not be confused with its much smaller exponential
correction.


\subsection{Experimental guidance and remaining questions}

The delivered filter verifier was replayed on an isolated copy with
mpmath 1.3.0 and SymPy 1.14.0 at seventy decimal digits. Its 165 exact
checks include finite operator inversions, conjugate-node interpolation
and the first nonlinear inverse coefficients. All 46 numerical checks
passed: regularized Stieltjes extraction against positive-integral jets,
tilted Borel integrals, noninteger completion, six neutral roots and the
small- and large-target inverse charts. A separate replay already filed
by the parallel intake is retained as independent run evidence.

For the new neutral inverse theorem, an independent implementation uses
Cauchy extraction of order jets from the entire Hurwitz realization,
with 64-to-96-point refinement, instead of tiny-step order
differentiation. It solves the neutral equation and the physical inverse
separately, retaining the moving-order term in
\eqref{ojf:eq:moving-chain-rule}. For $N=1$, centers $m=0,1,2$ and
$w=10,20$, the displacement agrees with $C_w/P'(w)$ at the predicted
leading scale. At $m=0,w=20$, their ratio is approximately
$1.00000000065748$. The logarithmic leading formula alone converges
more slowly: the exact derivative retains the lower powers of $\log w$.
These finite floating-point comparisons are diagnostics, not interval
certificates or proofs of the infinite assertions. The analytic proof
and exact coefficient identity are given above.

The higher-zero-jet Borel germs, conditioning of logarithmic-node
interpolation, growing filter size, and the exact complex radius of the
small-target inverse remain open directions. The neutral inverse theorem
settles fixed $m,N$ transport along the constructed curve, without
extending it to those different uniform or global questions.
